%!TEX root = main.tex

% \section{Overlapping-Induced Slowdown Prediction}
\section{Overlapping Slowdown Prediction} 
\label{sec:overlap_pred} 
% In \cref{subsec:challenges}, we analyzed the significant slowdown caused by overlapping kernels on a single GPU device. 
\noindent Accurately capturing and modeling the resource contention between overlapping kernels presents substantial challenges in predicting the slowdown factor as discussed in \cref{subsec:challenges}. 
To address this, \sys introduces an ML-based slowdown prediction model with a number of hardware- and software-specific features, which we present here.


% \subsection{ML-based Slowdown Prediction Model} 
% \label{subsec:slowdown_model}
A GPU devices executes concurrent computation and communication kernels with spatial multitasking, where each kernel is assigned to separate streams with a subset of the SMs.
At the same time they also contend for shared resources such as cache and DRAM bandwidth.
% Traditionally, a GPU device is reserved for a single stream, allowing it to monopolize all available resources. However, launching multiple streams concurrently results in spatial multitasking, where each stream is allocated a subset of streaming multiprocessors (SMs). For instance, in distributed DNN training, the stream dedicated to the all-reduce kernel may be allocated 16 SMs out of the 108 available on an NVIDIA A100 GPU. Additionally, shared resources such as cache and DRAM bandwidth are also distributed among the streams.
These factors results in performance interference, and \sys needs to capture them in predicting the slowdown. 
% To address the issue of kernel interference, we adopt a black-box approach using an ML model to learn the latent GPU scheduling mechanisms and predict the slowdown factor.

% \noindent\textbf{Key design.} We leverage machine learning techniques to model the interference between computation and communication kernels. Our ML model predicts the inverse of the slowdown factor, focusing on leveraging single-device training performance, which is typically easier to obtain. The model is trained using kernel performance data from single-device runs and the parameters of communication setting to predict their interaction and resulting slowdown.

\begin{table}[t] 
    \centering 
    \resizebox{\columnwidth}{!}{ 
        \begin{tabular}{@{}cll@{}} 
            \toprule \multicolumn{1}{l}{Type} & Feature & Specification \\ \midrule 
            \multirow{6}{*}{\begin{tabular}[c]{@{}c@{}}Comm. \\ details \end{tabular}}
            & Protocol & Communication protocol used within NCCL \\ 
            & Algorithm & Communication topology algorithm applied in NCCL \\ 
            & Collective & Name of the collective function in the kernel \\ 
            & Bucket size & Size of the tensor bucket used in collective communication \\ 
            & Channel number & Number of communication channels in NCCL \\ 
            \midrule \multirow{8}{*}{\begin{tabular}[c]{@{}@{}c@{}}Comp. \\ Kernel \\ Metric \end{tabular}} 
            & Running time &  Running time profiled on single-device training \\
            & Compute throughput & Average SM throughput\\ 
            & Memory throughput & Percentage of cycles where DRAM was active\\ 
            & DRAM throughput & Peak DRAM throughput \\ 
            & Achieved occupancy & Average percentage of active warps on each SM \\ 
            & L1 hit rate & Hit rate for L1 cache \\ 
            & L2 hit rate & Hit rate for L2 cache \\ 
            \bottomrule 
        \end{tabular}} 
    \caption{Features in the slowdown prediction model. } 
    \label{table:features_interference} 
    \vspace{-5mm}
\end{table}


\noindent\textbf{Feature extraction.} 
Table~\ref{table:features_interference} summarizes the key features used by our model. We categorize features into two groups: communication-related details and computation kernel metrics. 
For communication, we consider features such as protocol and algorithm employed by NCCL, CC kernel type, bucket size, and channel configuration. These parameters capture the complexity and overheads of communication: bucket size, for instance, determines when a communication kernel will be launched after accumulating one bucket of data. 
On the computation side, we profile kernels at a fine-grained level on a single device to measure their baseline performance without overlapping. 
We collect running times, compute throughput, memory utilization, and cache statistics using NVIDIA Nsight Systems and Nsight Compute. 
For instance, achieved occupancy and SM activity levels indicate compute intensity, while DRAM utilization and memory throughput---along with L1/L2 hit rates---shed light on memory access latency and potential bottlenecks. Additionally, we classify kernels by their functional types (e.g., matrix multiplication, layer normalization) to further refine our predictions.

\noindent\textbf{XGBoost-based model.} We leverage XGBoost~\cite{xgboost}, a tree-based gradient boosting framework, to predict the slowdown factors caused by overlapping kernels with the above features. 
XGBoost is well-suited to this task for two key reasons. First, our dataset contains heterogeneous features that are both numerical (e.g., SM throughput, DRAM utilization) and categorical (e.g., GPU type, kernel type), a scenario where tree-based models excel. 
Second, our target slowdown factors, influenced by the interplay of diverse system conditions, remain bounded within a range well-captured by the training set. As we expand our dataset---incorporating additional profiles from various models, GPU architectures, and cluster scales---the model's generalization ability improves, enabling accurate slowdown predictions across increasingly complex and heterogeneous environments.

% \noindent\textbf{Training data preparation.} We construct a comprehensive training dataset from two primary sources. First, we gather detailed Nsight profiling data from our GPU clusters, capturing a wide range of operational scenarios. Second, we conduct targeted experiments that deliberately vary model architectures, and GPU types to broaden our dataset's coverage. By continuously incorporating new profiles as our system evolves, we enable the model to remain accurate and robust, effectively predicting slowdown factors for a wide spectrum of compute and communication kernels.
